% SPDX-License-Identifier: CC-BY-4.0
\documentclass[11pt,reqno]{amsart}
\usepackage[T1]{fontenc}
\usepackage{lmodern,amssymb,mathtools,microtype}
\usepackage[margin=1.15in]{geometry}
\usepackage{xcolor}
\usepackage[colorlinks=true,linkcolor=blue!45!black,citecolor=blue!45!black,urlcolor=blue!45!black]{hyperref}
\hypersetup{pdftitle={A Solution to the Gerstenhaber Problem for Three Commuting Matrices},pdfauthor={Yan-cheng Lin}}
\newtheorem{theorem}{Theorem}[section]
\newtheorem{proposition}[theorem]{Proposition}
\newtheorem{lemma}[theorem]{Lemma}
\newtheorem{corollary}[theorem]{Corollary}
\theoremstyle{definition}
\newtheorem{example}[theorem]{Example}
\newcommand{\Mat}{\operatorname{Mat}}
\newcommand{\End}{\operatorname{End}}
\newcommand{\Hom}{\operatorname{Hom}}
\newcommand{\im}{\operatorname{im}}
\newcommand{\ord}{\operatorname{ord}}
\newcommand{\length}{\operatorname{length}}
\newcommand{\coker}{\operatorname{coker}}
\newcommand{\tr}{\operatorname{tr}}
\newcommand{\GL}{\operatorname{GL}}
\newcommand{\OO}{\mathcal O}
\newcommand{\A}{\mathbb A}
\newcommand{\CC}{\mathbb C}
\title[The Gerstenhaber Problem for Three Commuting Matrices]{A Solution to the Gerstenhaber Problem for Three Commuting Matrices}
\author{Yan-cheng Lin}
\address{Anhui University of Chinese Medicine, China (present affiliation)}
\email{lin522026298@outlook.com}
\date{November 18, 2026}
\subjclass[2020]{Primary 15A27; Secondary 13C05, 14M12}
\keywords{Commuting matrices, faithful modules, determinantal ideals, nilpotent centralizers}
\begin{document}

\begin{abstract}
Three pairwise commuting endomorphisms of an \(n\)-dimensional vector space
generate a unital algebra of dimension at most \(n\). The bound holds over
every field and is sharp in every dimension. The proof fixes the number of
Jordan blocks of one generator and allows their lengths to vary. A violation
of the bound would give a pole of a determinantal minor. Along its divisor
the second generator would act without a cyclic vector both on the
cokernel and on the kernel of the first, after the first has been shifted
to a nilpotent operator. These two conditions confine the image of the
divisor under a finite map to a set of codimension at least two.
A companion Lean formalization proves the theorem over fields of
characteristic zero.
\end{abstract}
\maketitle

\section{Introduction}
How small a vector space can carry a faithful representation of a given
commutative algebra? The regular representation of a finite-dimensional
algebra \(B\) has dimension \(\dim B\). For an algebra generated by two
elements this is the smallest possible dimension, by the commuting-pair
theorem of Gerstenhaber~\cite{Ger} and the earlier irreducibility theorem
of Motzkin and Taussky~\cite{MT}. Four generators already behave
differently. The four matrix units \(E_{13},E_{14},E_{23},E_{24}\) in
\(\Mat_4(K)\) have pairwise products zero. Together with the identity
they span a five-dimensional algebra. The case between these two is the
three-matrix Gerstenhaber problem.

\begin{theorem}\label{thm:main}
Let \(K\) be a field and let \(T,Y,Z\) be pairwise commuting endomorphisms
of an \(n\)-dimensional \(K\)-vector space. Then
\[
                 \dim_K K[T,Y,Z]\leq n .
\]
Here the algebra contains the identity.
\end{theorem}

A single nilpotent Jordan block of size \(n\), together with two zero
matrices, attains equality for every \(n\geq 1\).
Equivalently, every finite-dimensional commutative \(K\)-algebra generated
by at most three elements has smallest faithful representation dimension
equal to its dimension. This formulation is classical
\cite{Bergman,RS}. Rajchgot and Satriano~\cite{RS} established the
inequality for certain extensions of cyclic modules. Cherny, Satriano
and Song~\cite{CSS} and Cherny, Le Quang and Satriano~\cite{CLS} treated
combinatorial families. A previous note of the author~\cite{Lin} treated
complex triples whose algebra contains an element with eigenspaces of
dimension at most four. For an arbitrary triple of size \(n\geq2\),
Luo, Song and Zeng~\cite[Theorem~1.1]{LSZ} obtained \(2n-2\), following
a linear bound for each fixed number of generators~\cite{LLSZ}.

The two-matrix proof suggests deforming a commuting pair to one with a
regular member. For triples the variety of commuting matrices can be
reducible~\cite{Gur}. There is another space in which to deform the
problem. After primary decomposition take \(T\) nilpotent, and view the
underlying space as a module \(W\) over \(k[[t]]\), with \(t\) acting by
\(T\). The number \(r\) of Jordan blocks of \(T\) is the number of
generators of this module. Once \(r\) is fixed, all choices of block
lengths can be described by a single square presentation
\[
                        W=k[[t]]^r/Pk[[t]]^r.
\]
The other two generators lift to matrices on the free module. Keeping
their actions on its relation lattice as well records two views of the
same representation. At a minimal presentation these give actions on
\(W/tW\) and on \(\ker t\), respectively. A cyclic vector on either
side (using the dual on the kernel side) forces the desired bound.

The reason both sides matter is a dimension count. A failure of the
inequality can be detected by the order of a single minor. It persists
along arcs transverse to a divisor in a space of dimension \(3r^2\).
By choosing \(r\) least among all failures, one can apply the preceding
cyclic-vector test at every general point of this divisor, even though
the corresponding \(W\) can have a different length. The two necessary
failures of cyclicity leave room for at most \(3r^2-2\) parameters.
Sections~\ref{sec:lattice}--\ref{sec:count} make this argument over an
algebraically closed field. The passage to an arbitrary field is given
at the end of Section~\ref{sec:count}.

\section{A length formula for commuting actions}\label{sec:lattice}
Let \(k\) be a field, \(\OO=k[[t]]\), and \(v\) the \(t\)-adic order.
Put \(v(0)=+\infty\). Take \(P\in\Mat_r(\OO)\) with \(p=\det P\ne0\),
where \(r\geq1\). The module \(W=\coker P\) has \(k\)-dimension
\(n=v(p)\) by elementary divisors. An \(\OO\)-linear endomorphism of
\(W\) lifts to the free module by lifting the images of its standard
generators. For integral matrices \(Y,Z\) on the free module, form
\begin{equation}\label{eq:coords}
         Y'=P^{-1}YP,\qquad Z'=P^{-1}ZP,\qquad C=P^{-1}[Y,Z].
\end{equation}
Then \(Y,Z\) induce commuting endomorphisms \(\bar Y,\bar Z\) of \(W\)
precisely when all three matrices in \eqref{eq:coords} are integral.
Integrality of \(Y',Z'\) expresses stability of the relation lattice.
Integrality of \(C\) places the image of the commutator in that lattice.

Write \(\operatorname{vec}(U)\) for the column obtained by stacking the
columns of \(U\) in order. Form the matrix
\begin{equation}\label{eq:matrix}
  M=\left[\,P^{\oplus r}\ \middle|\
       \big(\operatorname{vec}(Y^iZ^j)\big)_{0\leq i,j<r}\,\right]
\end{equation}
with \(r^2\) rows. Let \(H\) be its ideal of maximal minors. In particular
\(p^r\in H\).
\begin{proposition}\label{prop:length}
Under the integrality conditions \eqref{eq:coords}, the image algebra
\(B=k[t,\bar Y,\bar Z]\subseteq\End_k(W)\) satisfies
\[
                      \dim_k B=r\,v(p)-v(H).
\]
The presentation may have redundant generators.
\end{proposition}
\begin{proof}
Evaluation on the \(r\) standard generators embeds \(\End_{\OO}(W)\)
as an \(\OO\)-submodule of \(W^r\). Denote by \(L\) the image of \(B\).
Every power series in \(t\) acts on \(W\) by a polynomial since \(t\)
is nilpotent there. Thus \(L\) is the \(\OO\)-span of evaluations of
words in \(\bar Y,\bar Z\). The endomorphisms commute on \(W\), and the
monic characteristic polynomials of their respective \(r\)-by-\(r\)
lifts annihilate them. Reducing the two exponents separately gives the
words in \eqref{eq:matrix}. Their evaluations are exactly the indicated
columns modulo \(P^{\oplus r}\). It follows that \(\coker M=W^r/L\).
Smith normal form identifies its length with \(v(H)\). Length over
\(\OO\) equals \(k\)-dimension for a finite-length module, and \(W^r\)
has length \(rn\).
\end{proof}

\begin{example}\label{ex:equal}
Take \(P=t^2I_2,\ Y=tE_{12},\ Z=tE_{21}\). Here \(n=4\), and \(B\)
has basis \(1,t,\bar Y,\bar Z\), since the last three elements are linearly
independent and all their products vanish. Both \(W/tW\) and \(\ker t\)
have dimension two, with zero action of \(\bar Y,\bar Z\). Neither \(W\)
nor its dual is cyclic. The lifts have divided commutator
\(C=E_{11}-E_{22}\), and the length formula gives \(v(H)=4\).
Any other lifts are congruent to these modulo \(t^2\), so their
commutator still equals \(t^2(E_{11}-E_{22})\) modulo \(t^3\).
Thus these commuting quotient actions have no commuting lifts in this
presentation.
\end{example}

We now let \(k\) be algebraically closed. For fixed \(r\), let \(G_r\)
be the reduced closure in \((\Mat_r)^6\) of the graph of
\eqref{eq:coords} on the open set \(\det P\ne0\) of triples \(P,Y,Z\).
The graph over this open set is \(\GL_r\times(\Mat_r)^2\). It follows
that \(G_r\) is irreducible of dimension \(3r^2\). Each minor in
\eqref{eq:matrix} is a regular function on \(G_r\). An integral
presentation gives an arc in \(G_r\), since its generic point is on the
graph. Conversely, any arc with \(p\ne0\) at its generic point satisfies
\(YP=PY'\), \(ZP=PZ'\), and \([Y,Z]=PC\). It therefore gives a commuting
triple on \(\coker P\).

\section{From one presentation to a divisor}\label{sec:divisor}
If \(\dim B>n\), Proposition~\ref{prop:length} singles out a maximal
minor \(D\) with \(v(D)<(r-1)v(p)\). The function \(D/p^{r-1}\) then
has negative order on one arc. We record the geometric consequence in
a form that applies to other length inequalities.

\begin{lemma}\label{lem:divisor}
Let \(X\) be an irreducible affine variety of dimension \(d\) over an
algebraically closed field \(k\). Let \(a,b\) be nonzero regular functions
and suppose \(X_b\) is normal. Let \(e\geq1\). If an arc
\(\gamma:\operatorname{Spec} k[[t]]\to X\) has
\(v(\gamma^*a)<e\,v(\gamma^*b)<\infty\), there is an irreducible closed
subvariety \(S\subset X\) of dimension \(d-1\) and a dense set of closed
points of \(S\) through each of which passes an arc satisfying the same
strict inequality. These arcs can be chosen with the same orders for
\(a\) and \(b\).
\end{lemma}
\begin{proof}
Write \(\nu:\widetilde X\to X\) for the normalization. It is finite
and an isomorphism over \(X_b\). The generic point of the arc lifts there,
and properness extends the lift over \(k[[t]]\). Hence \(a/b^e\) is not
regular on \(\widetilde X\), as its value would then belong to \(k[[t]]\).
A Noetherian normal domain is the intersection
of its height-one localizations in its fraction field. Some prime divisor
\(E\subset\widetilde X\) thus has
\(m=\ord_E a<e\,\ord_E b=e h\). Both orders are nonnegative and \(h>0\).
The finite map \(E\to S=\nu(E)\) preserves dimension, so \(\dim S=d-1\).
See \cite[Tags 0BXR and 031T]{Stacks} for the normalization and
intersection statements.

At the generic point of \(E\) the local ring of \(\widetilde X\) is a
discrete valuation ring. Shrinking about that point, take an equation
\(u\) for \(E\) with \(a/u^m\) and \(b/u^h\) units. The field \(k\) is
perfect. Both \(\widetilde X\) and the reduced divisor \(E\) are smooth
along a dense open of \(E\), using normality in codimension one for the
ambient variety. At any closed point of this open, complete \(u\) to
regular parameters \(u,w_2,\ldots,w_d\). The completed local ring is
\(k[[u,w_2,\ldots,w_d]]\) with the given copy of \(k\)
\cite[Tag 0C0S]{Stacks}. Set \(u=t\) and all \(w_i=0\). The two unit
factors stay units, so the orders on this arc are \(m,h\). Its image in
\(X\) has the same orders. The images of the chosen closed points are
dense in \(S\).
\end{proof}

\section{The two actions at a closed point}\label{sec:ends}
Suppose the theorem fails over the fixed algebraically closed field \(k\).
A simultaneous primary decomposition splits the space according to the
joint eigenvalues of the three operators. Restriction embeds the image
algebra into the product of the image algebras on the summands. Thus a
failure persists on some summand. Subtracting its eigenvalue from \(T\),
we obtain a failure with \(T\) nilpotent.

From now on let \(r\) be the smallest number of Jordan blocks of \(T\)
among \emph{all} such failures over \(k\). If \(r=1\), a cyclic vector
for \(T\) identifies \(W\) with \(k[t]/(t^n)\), on which every operator commuting
with \(T\) is multiplication by an element of that algebra. Hence \(r\geq2\).
Every failing triple with nilpotent \(T\) and at most \(r\) blocks has
exactly \(r\) blocks and a single joint character. Indeed the block counts
add over joint primary summands, each nonzero summand contributes at least
one block, and some summand would still fail.

Present one such \(W\) on \(r\) generators over \(\OO\), using one cyclic
summand for each block. Apply Lemma~\ref{lem:divisor} to \(G_r\) with
\(a=D,\ b=p,\ e=r-1\), using smoothness of the open graph. At the resulting dense
set of closed points of a \((3r^2-1)\)-dimensional subvariety, the
transverse arc gives a module of dimension \(v(p)\) with at most \(r\)
generators. The order inequality and Proposition~\ref{prop:length} make
this module a failure even if a different minor now has least order.
Its block count is consequently \(r\). Since \(W/tW=\coker P(0)\),
this says \(P(0)=0\).

\begin{lemma}\label{lem:ends}
For a presentation satisfying \eqref{eq:coords} with \(P(0)=0\), the
actions of \(\bar Y,\bar Z\) on \(W/tW\) are given by \(Y(0),Z(0)\).
On \(\ker(t:W\to W)\) they are given in a suitable basis by
\(Y'(0),Z'(0)\). If \(Y(0)\) or \(Y'(0)\) has a cyclic vector, then
\(\dim_k B=\dim_k W\).
\end{lemma}
\begin{proof}
Reduction modulo \(t\) gives the first assertion. Write \(P=tQ\).
The map \(k^r\to\ker t\) sending a vector with lift \(v\) to
\(Qv\bmod P\OO^r\) is independent of the lift. If \(Qv=Pw=tQw\),
injectivity of \(Q\) gives \(v=tw\). Conversely \(tw=Pv\) gives
\(w=Qv\), proving surjectivity. The identities \(YQ=QY'\) and \(ZQ=QZ'\)
identify the two actions on the kernel.

If \(Y(0)\) has a cyclic vector, lift it to \(w\in W\). The quotient
\(W/Bw\) equals \(t(W/Bw)\), hence vanishes by nilpotence. Evaluation
\(B\to W,\ b\mapsto bw\) is surjective. For \(b\) in its kernel, commutativity gives
\(bBw=0\), and faithfulness gives \(b=0\). For the second case pass to
\(W^\vee=\Hom_k(W,k)\). Its quotient by \(t\) identifies by restriction
with \((\ker t)^\vee\). A cyclic matrix has cyclic transpose, and the
algebra on \(W^\vee\) has the same dimension as \(B\).
\end{proof}

Let \((\lambda,\mu)\) be the joint eigenvalues of \(Y,Z\) on one of these
modules. Both end pairs inherit these same eigenvalues. After subtracting
\((\lambda I,\mu I)\), each is a pair of commuting nilpotents. In each
pair the first nilpotent has at least two Jordan blocks, since a single
block has a cyclic vector. The pair of scalars may change from one point to
the next. It is shared by the two ends at each point.

\section{Dimension of the possible centers}\label{sec:count}
\begin{lemma}\label{lem:stratum}
Fix a partition of \(r\) with \(s\) parts. The variety of pairs
\((U,V)\in\Mat_r(k)^2\) with \(U\) nilpotent of this Jordan type and
\(V\) nilpotent commuting with \(U\) has dimension \(r^2-s\).
\end{lemma}
\begin{proof}
Regard \(k^r\) as \(M=\bigoplus_h (k[u]/u^h)^{m_h}\), with \(u\)
acting by \(U\). Let \(E=\End_{k[u]}(M)\) and \(c=\dim E\). Its units
form an open set, so the conjugacy orbit of \(U\) has dimension \(r^2-c\).

To count the fiber over \(U\), consider the projection to the
constant terms of the diagonal blocks of each fixed length,
\[
                E\longrightarrow\prod_h \Mat_{m_h}(k).
\]
This is a surjective algebra map. In composing blocks, any passage from
one length to another and back acquires a factor of \(u\) on returning.
The kernel is nilpotent. Indeed on \(M/uM\) every map from a shorter
summand to a longer one vanishes, and a kernel element has zero
equal-length diagonal blocks. Products of as many kernel elements as
there are distinct lengths therefore send \(M\) into \(uM\). Iterating
sends \(M\) to zero. Consequently the nilpotents in \(E\) are exactly
the inverse image of the product of the nilpotent cones. In size \(m\)
this cone has dimension \(m^2-m\), since the centralizers in its finitely
many Jordan orbits have dimension at least \(m\), with equality for one block.
(The block-diagonal polynomials in each Jordan block already give
\(m\) centralizing directions.) The dimension of the fiber is therefore
\(c-\sum_h m_h=c-s\). Adding the orbit dimension gives the formula.
\end{proof}

\begin{proof}[Proof of Theorem~\ref{thm:main}]
First work over the fixed algebraically closed field \(k\). Let
\(\mathcal N_{\geq2}\) consist of the pairs in Lemma~\ref{lem:stratum}
whose first member has at least two blocks, allowing all such partitions.
It has dimension at most \(r^2-2\). The closed points obtained in
Section~\ref{sec:ends} belong to the image of
\[
 \A^2\times\mathcal N_{\geq2}\times\mathcal N_{\geq2}\times\Mat_r
             \longrightarrow (\Mat_r)^6 .
\]
The map sets \(P=0\), adds the same two scalars to the two pairs, and
takes its last argument as \(C\). In particular the images of the two
pairs may satisfy further relations. Even allowing the two pairs and
\(C\) to vary throughout this product gives an image of dimension at most
\[
                         2+2(r^2-2)+r^2=3r^2-2 .
\]
Its closure contains the dense set of points from
Lemma~\ref{lem:divisor}, whose closure has dimension \(3r^2-1\). This
contradiction proves the theorem over \(k\).

For any field \(K\), extend a given triple to \(k=\overline K\).
Tensoring \(B=K[T,Y,Z]\hookrightarrow\Mat_n(K)\) with \(k\) is injective,
and its image is \(k[T,Y,Z]\), since the same words span both algebras.
Dimensions agree and the inequality descends. This also handles finite
and imperfect fields. In positive characteristic the ordered words in
\eqref{eq:matrix} use separate monic reductions, and the arcs of
Lemma~\ref{lem:divisor} use smooth closed points over \(\overline K\).
Both operations work for every block size.
\end{proof}

\section{Consequences and further questions}\label{sec:consequences}
The proof also yields divisibility statements on a single affine variety,
independent of the lengths of the Jordan blocks. Write \(R_r=k[G_r]\)
and \(R_r^\nu\) for its normalization, with \(k\) algebraically closed.
Let \(H_r\subset R_r\) be the ideal of \emph{all} maximal minors of
\eqref{eq:matrix}. For each \(r\geq1\) the following ideal is defined on
the normalization.
\begin{corollary}\label{cor:ideal}
One has \(H_r R_r^\nu\subseteq p^{r-1}R_r^\nu\).
Put \(J_r=p^{1-r}H_r R_r^\nu\). For any arc of \(G_r\) with generic point
in \(p\ne0\), pull \(J_r\) back along its lift to the normalization.
Then
\[
            v(J_r)=\dim_k W-\dim_k k[t,\bar Y,\bar Z].
\]
In particular the representations attaining equality are precisely the
arcs whose lifted centers lie outside \(V(J_r)\).
\end{corollary}
\begin{proof}
Take a nonzero minor \(D\). If \(D/p^{r-1}\) had a pole along a divisor
of the normalization, the closed-point construction in
Lemma~\ref{lem:divisor} would give an arc with \(v(D)<(r-1)v(p)\),
contradicting the theorem and Proposition~\ref{prop:length}. Normality
now gives the inclusion. The same proposition gives
\(v(H_r)-(r-1)v(p)=n-\dim_k B\), which proves the remaining assertions.
\end{proof}
This suggests a question about the original presentation space. For
which \(r\) does \(H_r\subseteq p^{r-1}R_r\) already hold in the
coordinate ring of the graph? Divisibility there would express every
normalized minor as a polynomial in the six matrix coordinates. The
corollary gives regularity after normalization.

There is also a quantitative consequence for faithful modules. Write
\(\mu_K(B)\) for the least dimension of a faithful module for a nonzero
finite-dimensional commutative algebra \(B\). Following~\cite{LSZ}, let
\(g_K(B)\) be the maximum rank of the bilinear forms
\((x,y)\mapsto\lambda(xy)\) for \(\lambda\in B^\vee\), and put
\(\widetilde g_K(B)=\lim_{q\to\infty}g_K(B^{\otimes q})^{1/q}\).
\begin{corollary}\label{cor:rank}
Over an infinite field \(K\),
\[
     \sup_{B}\frac{\mu_K(B)}{\widetilde g_K(B)}=2 ,
\]
where \(B\) ranges over the nonzero finite-dimensional commutative
\(K\)-algebras generated by at most three elements.
\end{corollary}
\begin{proof}
The regular module and Theorem~\ref{thm:main} give
\(\mu_K(B)=\dim_K B\). The supremum with this numerator is
\(\Gamma(3)/\Gamma(2)^2=2\) by \cite[Theorem~1.3]{LSZ}.
\end{proof}

\subsection*{A relation for four generators}
Write \(c_4(K)\) for the supremum of \(\dim_K B/n\) over algebras
generated by four commuting \(n\)-by-\(n\) matrices. Constructions in
\cite{LLSZ} and the upper bound in \cite[Corollary~1.4]{LSZ} give
\(2\leq c_4(K)\leq32/(3\pi)\). The fourth generator changes the lattice
calculation in a specific way. With three companions \(Y_1,Y_2,Y_3\)
of \(t\) one must keep three divided commutators \(C_{ij}\).
They satisfy a linear relation once the end actions are fixed.

\begin{proposition}\label{prop:jacobi}
Let \(Y_iP=PY'_i\) and \([Y_i,Y_j]=PC_{ij}\) over \(k[[t]]\), where
\(\det P\ne0\), \(1\leq i,j\leq3\), and \(C_{ji}=-C_{ij}\).
At a center with \(P(0)=0\), write \(A_i=Y_i(0)\), \(A'_i=Y'_i(0)\),
and \(D_{ij}=C_{ij}(0)\). Then
\[
              \sum_{\rm cyc}\big(A'_iD_{jk}-D_{jk}A_i\big)=0 .
\]
For fixed end actions the space of triples of matrices satisfying this
linear equation has dimension
\[
 2r^2+\dim_k\{S\in\Mat_r(k): S A'_i=A_i S\text{ for }1\leq i\leq3\}.
\]
\end{proposition}
\begin{proof}
Insert the expressions \(PC_{jk}\) into
\(\sum_{\rm cyc}[Y_i,[Y_j,Y_k]]=0\), move \(P\) to the left using
\(Y_iP=PY'_i\), and cancel \(P\). Evaluation at zero gives the relation.
Pair the target matrix space with itself by \((S,M)\mapsto\tr(SM)\).
The annihilator of the image consists of the simultaneous intertwiners
in the display, by cyclic invariance and nondegeneracy of the trace
pairing. Rank--nullity gives the formula.
\end{proof}

How large is the locus of divided commutators actually attained for
prescribed end actions in this linear space? A description would constrain
the third piece of the four-generator parameter count. The three-generator
argument treats its single commutator coordinate as entirely free.

\appendix
\section{The characteristic-zero formalization}\label{sec:lean}
The companion proves the bound for three named matrices in each of the
following settings.
\begin{center}
\begin{tabular}{ll}
complex coefficients & \texttt{G3Assembly.complex\_three\_matrices} \\
characteristic zero & \texttt{G3Scalar.charZero\_three\_matrices}.
\end{tabular}
\end{center}
Both concern the unital subalgebra generated by the three matrices under
their pairwise commutation hypotheses.
The positive-characteristic case in Theorem~\ref{thm:main} is proved by
the closed-point argument in the text. The source map and reproduction
instructions accompany the paper.

There is one difference of mathematical interest between the two proofs.
The formal proof works at the generic point of a divisor. It adjoins
\(q=p^{r-1}/D\) to the graph ring. On the original arc this reciprocal
has positive order, so it is a nonzero nonunit upstairs in the enlarged
ring (the specialization is taken in \(k((t))\)). A prime minimal over
\(q\) in the finite normalization gives a discrete valuation ring \(V\).
In the residue field \(q=0\), so the six matrix coordinates generate a
field of transcendence degree \(3r^2-1\) over the base, up to a finite
extension. Thus the dimension to be contradicted is still visible in
the original coordinates.

To realize this generic data by a finite-dimensional representation, take
\(N>\max(v(D),v(p))\). In characteristic zero the truncated discrete
valuation ring has a coefficient field compatible with the ground field.
In power series over the residue field, lift the six matrices from this
jet. The three relation errors are divisible by \(t^N\). Multiplication
by \(P^{-1}=\operatorname{adj}(P)/p\) repairs the last three matrices
without changing their reductions. The orders of the two determinants
persist. The least block count is now taken across algebraically closed
fields of characteristic zero, which permits this change of residue
field. The rest is the count on the two ends, expressed as a bound on
the number of algebraically independent entries. This accounts for the
two routes from a single minor to the same dimension defect.

\subsection*{Acknowledgment}
GPT-series agents assisted with the mathematical exploration, proof
development, Lean formalization, and preparation of the manuscript.
The author is responsible for the mathematical statements and the
accompanying proof materials released with this work.

\begin{thebibliography}{99}
\bibitem{Bergman}
G.~M. Bergman, \emph{Commuting matrices, and modules over Artinian local
rings}, \href{https://arxiv.org/abs/1309.0053}{arXiv:1309.0053} (2013).
\bibitem{CLS}
R. Cherny, T.~A. Le Quang and M. Satriano, \emph{On combinatorial algebras
generated by three commuting matrices},
\href{https://arxiv.org/abs/2511.22092}{arXiv:2511.22092} (2025).
\bibitem{CSS}
R. Cherny, M. Satriano and Y. Song, \emph{On the algebra generated by three
commuting matrices: combinatorial cases}, Electron. J. Combin.
\textbf{31} (2024), no.~4, Paper 4.43.
\href{https://doi.org/10.37236/12909}{doi:10.37236/12909}.
\bibitem{Ger}
M. Gerstenhaber, \emph{On dominance and varieties of commuting matrices},
Ann. of Math. (2) \textbf{73} (1961), 324--348.
\href{https://doi.org/10.2307/1970336}{doi:10.2307/1970336}.
\bibitem{Gur}
R. Guralnick, \emph{A note on commuting pairs of matrices}, Linear
Multilinear Algebra \textbf{31} (1992), 71--75.
\href{https://doi.org/10.1080/03081089208818123}{doi:10.1080/03081089208818123}.
\bibitem{LLSZ}
J. Li, C. Luo, X. Song and Y. Zeng, \emph{Linear bounds and sharp
exponential growth for commuting matrix algebra},
\href{https://arxiv.org/abs/2611.08135}{arXiv:2611.08135} (2026).
\bibitem{Lin}
Y.-c. Lin, \emph{Three-generated matrix algebras containing a four-regular
element}, Zenodo preprint (2026).
\href{https://doi.org/10.5281/zenodo.24242097}{doi:10.5281/zenodo.24242097}.
\bibitem{LSZ}
C. Luo, X. Song and Y. Zeng, \emph{Linear bound for commuting matrix algebra
II}, \href{https://arxiv.org/abs/2611.11550}{arXiv:2611.11550} (2026).
\bibitem{MT}
T.~S. Motzkin and O. Taussky, \emph{Pairs of matrices with property \(L\).
II}, Trans. Amer. Math. Soc. \textbf{80} (1955), 387--401.
\href{https://doi.org/10.1090/S0002-9947-1955-0086781-5}{doi:10.1090/S0002-9947-1955-0086781-5}.
\bibitem{RS}
J. Rajchgot and M. Satriano, \emph{New classes of examples satisfying the
three matrix analog of Gerstenhaber's theorem}, J. Algebra \textbf{516}
(2018), 245--270.
\href{https://doi.org/10.1016/j.jalgebra.2018.09.020}{doi:10.1016/j.jalgebra.2018.09.020}.
\bibitem{Stacks}
The Stacks Project Authors, \emph{The Stacks Project},
Tags \href{https://stacks.math.columbia.edu/tag/0BXR}{0BXR},
\href{https://stacks.math.columbia.edu/tag/031T}{031T} and
\href{https://stacks.math.columbia.edu/tag/0C0S}{0C0S}.
\end{thebibliography}
\end{document}
